\section{Expectation and Variance Algebra}

\subsection{Expectation}

Let \(\Omega\) be finite or uncountable, and let \(\functiondomain{X}{\Omega}{\RR}\) be a discrete random variable.

We say \(X\) is \emph{non-negative} if \(X \geq 0\). We first define the expectation for a non-negative random variable (see later for motivation).

\begin{definition}[Expectation of Non-Negative Random Variable]
    \label{def:expectation-non-negative}%
    The \emph{expectation}, or \emph{mean/average} of a non-negative random variable is given by
    \[
        \Expt \bqty{X} = \sum_{\omega \in \Omega} \Prob\pqty{\Bqty{\omega}} \cdot X\pqty{\omega}.
    \]
\end{definition}

Equivalently, we may define the expectation as follows. If we let the range of the random variable be \(\Omega_X = \set{X\pqty{x}}{\omega \in \Omega}\), and partition the sample space as
\[
    \Omega = \bigcup_{x \in \Omega_X} \Bqty{X = x},
\]
then the expectation satisfies
\begin{align*}
    \Expt\bqty{X} &= \sum_{\omega \in \Omega} X\pqty{\omega} \Prob\pqty{\Bqty{\omega}} \\
    &= \sum_{x \in \Omega_X} \sum_{\omega \in \Bqty{X = x}} X\pqty{\omega} \cdot \Prob\pqty{\Bqty{\omega}} \\
    &= \sum_{x \in \Omega_X} \sum_{\omega \in \Bqty{X = x}} x \cdot \Prob\pqty{\Bqty{\omega}} \\
    &= \sum_{x \in \Omega_X} x \sum_{\omega \in \Bqty{X = x}} \cdot \Prob\pqty{\Bqty{\omega}} \\
    &= \sum_{x \in \Omega_X} x \cdot \Prob\pqty{X = x}.
\end{align*}

In other words, the expectation of \(X\) is the average of the values \(x\) taken by \(X\), with weights given by the corresponding probabilities \(\Prob\pqty{X = x}\).

\begin{example}[Expectation of Binomial Random Variable]
    \label{ex:expectation-binomial}%
    Let \(X \sim \Binomial\pqty{N, p}\). This has probability mass function
    \[
        \Prob\pqty{X = k} = \binom{N}{k} p^k \pqty{1 - p}^{N - k}.
    \]

    By definition,
    \begin{align*}
        \Expt\bqty{X} &= \sum_{k = 0}^{N} k \Prob\pqty{X = k} \\
        &= \sum_{k = 0}^{N} k \binom{N}{k} p^k \pqty{1 - p}^{N - k} \\
        &= \sum_{k = 0}^{N} k \cdot \frac{N!}{\pqty{N - k}! k!} p^k \pqty{1 - p}^{N - k} \\
        &= \sum_{k = 1}^{N} N \cdot \frac{\pqty{N - 1}!}{\pqty{N - k}! \pqty{k - 1}!} p^k \pqty{1 - p}^{N - k} \\
        &= \sum_{k = 0}^{N - 1} N \binom{N - 1}{k} p^{k} \pqty{1 - p}^{N - k - 1} p \\
        &= Np \pqty{p + 1 - p}^{N - 1} \\
        &= Np.
    \end{align*}
\end{example}

\begin{example}[Expectation of Poisson Random Variable]
    \label{ex:expectation-poisson}%
    Let \(X \sim \Poisson\pqty{\lambda}\) with \(\lambda > 0\). We have
    \[
        \Prob\pqty{X = k} = e^{-\lambda} \cdot \frac{\lambda^k}{k!}.
    \]

    By definition,
    \begin{align*}
        \Expt\bqty{X} &= \sum_{k = 0}^{\infty} k \cdot \Prob\pqty{X = k} \\
        &= \sum_{k = 1}^{\infty} k \cdot e^{-\lambda} \cdot \frac{\lambda^k}{k!} \\
        &= e^{-\lambda} \cdot \lambda \cdot \sum_{k = 1}^{\infty} \frac{\lambda^{k - 1}}{\pqty{k - 1}!} \\
        &=\lambda.
    \end{align*}
\end{example}

Now we define the expectation of a general random variable \(X\). The reason why we defined it for a non-negative random variable first, is due the sum being over a \emph{set}, which has no built-in ordering to specify the sum. It is always fine to rearrange a non-negative sum, since it either absolutely converges or always diverges, and rearranging does not change its sum.

However, if we straight-away define it for a general random variable using the same sum, and if the series is only \emph{conditionally convergent}, then changing the order of summation will change the limit, which makes the expectation change. Therefore, we will formalise this considering the sum of the positive and negative values separately.

Let \(X^+ = \max\Bqty{X, 0}\) and \(X^- = \min\Bqty{-X, 0}\).

We therefore have \(X = X^+ - X^-\) and \(\abs{X} = X^+ + X^-\).

Suppose \(X\) is discrete. Using \autoref{def:expectation-non-negative}, we have already defined \(\Expt\bqty{X^+}\) and \(\Expt\bqty{X^-}\).

\begin{definition}[Expectation of General Random Variable]
    \label{def:expectation-general-random-variable}%
    Let \(X\) be a discrete random variable and \(X^+\) and \(X^-\) be defined as above. If at least one of \(\Expt\bqty{X^+}\) and \(\Expt\bqty{X^-}\) is finite, then we define the \emph{expectation} of \(X\) as
    \[
        \Expt\bqty{X} = \Expt\bqty{X^+} - \Expt\bqty{X^-}.
    \]

    If both of them are infinite, then the expectation of \(X\) is \emph{undefined}.
\end{definition}

When we write \(\Expt\bqty{X}\), we implicitly imply that it is defined.

\begin{definition}[Integrablility of Random Variable]
    \label{def:integrability-random-variable}%
    If \(\Expt\bqty{\abs{X}} < \infty\), then we say \(X\) is \emph{integrable}.
\end{definition}

When \(\Expt\bqty{X}\) is well-defined, then we have
\[
    \Expt\bqty{X} = \sum_{x \in \Omega_X} x \Prob\pqty{X = x}.
\]

\begin{proposition}[Properties of Expectation]
    \label{prop:properties-expectation}%
    \begin{enumerate}
        \item If \(X \geq 0\), then \(\Expt\bqty{X} \geq 0\).
        \item If \(X \geq 0\) and \(\Expt\bqty{X} = 0 \), then \(\Prob\pqty{X = 0} = 1\).
        \item If \(c \in \RR\), then \(\Expt\bqty{cX} = c \Expt\bqty{X}\), and \(\Expt\bqty{c + X} = c + \Expt\bqty{X}\).
        \item If \(X\) and \(Y\) are random variables, then \(\Expt\bqty{X + Y} = \Expt\bqty{X} + \Expt\bqty{Y}\),
        \item Let \(c_1, \ldots, c_n \in \RR\), and \(X_1, \ldots, X_n\) be integrable random variables. Then we have
        \[
            \Expt \bqty{\sum_{i = 1}^{n} c_i X_i} = \sum_{i = 1}^{n} c_i \Expt \bqty{X_i}.
        \]
        \item Suppose \(X_1, X_2, \ldots\) are non-negative random variables. Then
        \[
            \Expt\bqty{\sum_{n \in \NN} X_n} = \sum_{n \in \NN} \Expt\bqty{X_n}.
        \]
    \end{enumerate}
\end{proposition}

\begin{proof}[on Linearity over Countable Collection]
    We have
    \begin{align*}
        \Expt\bqty{\sum_{n \in \NN} X_n} &= \sum_{\omega \in \Omega} \pqty{\sum_{n \in \NN} X_n \pqty{\omega}} \cdot \Prob\pqty{\Bqty{\omega}} \\
        &= \sum_{n \in \NN} \sum_{\omega \in \Omega} X_n \pqty{\omega} \Prob\pqty{\Bqty{\omega}} \\
        &= \sum_{n \in \NN} \Expt\bqty{X_n}.
    \end{align*}
\end{proof}

\begin{example}[Expectation of Indicator]
    \label{ex:expectation-indicator}%
    Recall the definition of the \emph{indicator} of an event \(A \in \eventSpace\) is \(\indicator\pqty{A}\). Let \(X = \indicator\pqty{A}\), then \(X\pqty{\omega} = \indicator\pqty{\omega \in A}\). Therefore, \(\Expt\bqty{X} = \Prob\pqty{A}\).
\end{example}

\begin{proposition}[Expectation of Function of Random Variable]
    \label{prop:expectation-function-random-variable}%
    Let \(\functiondomain{X}{\Omega}{\RR}\) be a random variable, and let \(\functiondomain{g}{\RR}{\RR}\). We define
    \[
        g\pqty{X} \pqty{\omega} = g\pqty{X\pqty{\omega}}.
    \]

    Then
    \[
        \Expt\bqty{g\pqty{X}} = \sum_{x \in \Omega_X} g\pqty{x} \cdot \Prob\pqty{X = x}.
    \]
\end{proposition}

\begin{proof}
    Set \(Y = g\pqty{X}\). Then we have
    \[
        \Expt\bqty{Y} = \sum_{y \in \Omega_Y} y \cdot \Prob\pqty{Y = y}.
    \]

    Now, note that the event can be simplified as
    \begin{align*}
        \Bqty{Y = y} &= \set{\omega \in \Omega}{Y \pqty{\omega} = y} \\
        &= \set{\omega \in \Omega}{g\pqty{X\pqty{\omega}} = y} \\
        &= \set{\omega \in \Omega}{X\pqty{\omega} \in g\inverse \pqty{\Bqty{y}}} \\
        &= \Bqty{X \in g\inverse \pqty{\Bqty{y}}}.
    \end{align*}

    Therefore,
    \begin{align*}
        \Expt\bqty{Y} &= \sum_{y \in \Omega_Y} y \cdot \Prob\pqty{X \in g\inverse\pqty{\Bqty{y}}} \\
        &= \sum_{y \in \Omega_Y} y \cdot \sum_{x \in g\inverse\pqty{\Bqty{y}}} \Prob\pqty{X = x} \\
        &= \sum_{y \in \Omega_Y} \sum_{x \in g\inverse\pqty{\Bqty{y}}} g\pqty{x} \cdot \Prob\pqty{X = x} \\
        &= \sum_{x \in \Omega_X} g\pqty{x} \Prob\pqty{X = x}.
    \end{align*}
\end{proof}

\begin{proposition}[Expectation of Natural Number-Values Random Variables]
    \label{prop:expt-natural-number-random-variable}%
    Suppose \(X \geq 0\) and takes integer values. Then we have
    \[
        \Expt\bqty{X} = \sum_{k = 1}^{\infty} \Prob\pqty{X \geq k} = \sum_{k = 0}^{\infty} \Prob\pqty{X > k}.
    \]
\end{proposition}

\begin{proof}
    Since \(X \geq 0\) and \(X \in \NN\), we have
    \[
        X = \sum_{k = 1}^{\infty} \indicator \pqty{X \geq k} = \sum_{k = 0}^{\infty} \indicator \pqty{X > k}.
    \]

    Indeed, this is true for any \(X = x \in \NN\), since
    \[
        x = \sum_{k = 1}^{x} 1 = \sum_{k = 1}^{x} \indicator \pqty{x \geq k} = \sum_{k = 1}^{\infty} \indicator \pqty{x \geq k}
    \]
    and similarly for the strict inequality case.

    Hence,
    \[
        X \pqty{\omega} = \sum_{k = 1}^{\infty} \indicator \pqty{X \pqty{\omega} \geq k} = \sum_{k = 0}^{\infty} \indicator \pqty{X \pqty{\omega} > k}.
    \]

    Taking the expectation, we get
    \begin{align*}
        \Expt\bqty{X} &= \Expt\bqty{\sum_{k = 1}^{\infty} \indicator\pqty{X \geq k}} \\
        &= \sum_{k = 1}^{\infty} \Expt\bqty{\indicator\pqty{X \geq k}} \\
        &= \sum_{k = 1}^{\infty} \Prob\pqty{X \geq k},
    \end{align*}
    and similarly for the strict case.
\end{proof}

The indicator random variable gives an alternative proof of \autoref{thm:inclusion-exclusion} (Inclusion--Exclusion Formula).

\begin{proof}[of \autoref{thm:inclusion-exclusion} (Inclusion--Exclusion Formula), using Indicator Functions]
    First, note that for events \(A, B \in \eventSpace\),
    \begin{itemize}
        \item \(\indicator\pqty{A\compl} = 1 - \indicator\pqty{A}\);
        \item \(\indicator\pqty{A \cap B} = \indicator\pqty{A} \cdot \indicator\pqty{B}\); and for the union,
        \item \(\indicator\pqty{A \cup B} = 1 - \pqty{1 - \indicator\pqty{A}} \pqty{1 - \indicator\pqty{B}}\).
    \end{itemize}

    More generally, for any \(A_1, A_2, \ldots, A_n \in \eventSpace\), we have
    \begin{align*}
        \indicator\pqty{A_1 \cup A_2 \cup \cdots \cup A_n} &= 1 - \prod_{i = 1}^{n} \pqty{1 - \indicator\pqty{A_i}} \\
        &= \sum_{i = 1}^{n} \indicator \pqty{A_i} - \sum_{i_1 < i_2} \indicator\pqty{A_{i_1} \cap A_{i_2}} + \cdots + \pqty{-1}^{n + 1} \indicator\pqty{A_1 \cap \cdots \cap A_n}.
    \end{align*}

    We may take the expectation on both sides, and use that
    \[
        \Expt\bqty{\indicator\pqty{A}} = \Prob\pqty{A}.
    \]

    Hence,
    \[
        \Prob\pqty{A_1 \cup A_2 \cup \cdots \cup A_n} = \sum_{k = 1}^{n} \sum_{1 \leq i_1 < i_2 < \cdots < i_k \leq n} \Prob\pqty{A_{i_1} \cap A_{i_2} \cap \cdots \cap A_{i_k}}.
    \]
\end{proof}

\subsection{Variance}

Defining the expectation allows us to define several other important things about a random variable.

\begin{definition}[Moment]
    \label{def:moment}%
    Let \(X\) be a random variable, and \(r \in \NN\). \(\Expt\bqty{X^r}\), i well-defined, is the \emph{\(r\)-th moment} of \(X\).
\end{definition}

\begin{definition}[Variance]
    \label{def:variance}%
    The \emph{variance} of \(X\) is given as
    \[
        \Var\pqty{X} = \Expt\bqty{\pqty{X - \Expt\bqty{X}}^2}.
    \]
\end{definition}

The variance is a `measure' of how well concentrated \(X\) is around its expectation. The smaller the variance, the more concentrated \(X\) is.

\begin{definition}[Standard Deviation]
    \label{def:std-deviation}%
    \(\sqrt{\Var\pqty{X}}\) is known as the \emph{standard deviation}.
\end{definition}

\begin{proposition}[Properties of the Variance]
    \label{prop:properties-of-variance}%
    \begin{itemize}
        \item \(\Var\pqty{X} \geq 0\).
        \item If \(\Var\pqty{X} = 0\), then \(\Prob\pqty{X = \Expt\bqty{X}} = 1\).
        \item If \(c \in \RR\), then \(\Var\pqty{cX} = c^2 \Var\pqty{X}\), and \(\Var\pqty{X + c} = \Var\pqty{X}\).
    \end{itemize}
\end{proposition}

\begin{proposition}[Alternative Formulation of Variance]
    \label{prop:alt-formulation}%
    We have
    \[
        \Var\pqty{X} = \Expt\bqty{X^2} - \Expt\bqty{X}^2.
    \]
\end{proposition}

\begin{proof}
    We have
    \begin{align*}
        \Var\pqty{X} &= \Expt\bqty{X^2 - 2X \Expt\bqty{X} + \pqty{\Expt\bqty{X}}^2} \\
        &= \Expt\bqty{X^2} - 2 \Expt\bqty{X} \Expt\bqty{X} + \pqty{\Expt\bqty{X}}^2 \\
        &= \Expt\bqty{X^2} - \pqty{\Expt\bqty{X}}^2.
    \end{align*}
\end{proof}

\begin{proposition}[Variance as a Minimalisation]
    \label{prop:variance-minimalisation}%
    We have
    \[
        \Var\pqty{X} = \min_{c \in \RR} \Expt\bqty{\pqty{X - c}^2},
    \]
    obtained for \(c = \Expt\bqty{X}\).
\end{proposition}

\begin{proof}
    We let
    \[
        f\pqty{c} = \Expt\bqty{\pqty{X - c}^2}.
    \]

    Then, we have
    \[
        f\pqty{c} = \Expt\bqty{X^2} = 2c \Expt\bqty{X} + c^2.
    \]

    Hence, differentiating gives that
    \[
        f'\pqty{c} = -2 \Expt\bqty{X} + 2c
    \]
    and hence
    \[
        \min_{c \in \RR} f\pqty{c} = f\pqty{\Expt\bqty{X}} = \Var\pqty{X}.
    \]
\end{proof}

\begin{examples}[Variances of Common Random Variables]
    \label{ex:variance-rv}%
    \begin{enumerate}
        \item If \(X \sim \Binomial\pqty{n, p}\), then \(\Expt\bqty{X} = np\).
        
        We have
        \[
            \Var\pqty{X} = \Expt\bqty{X\pqty{X - 1}} + \Expt\bqty{X} - \pqty{\Expt\bqty{X}}^2.
        \]

        We can find
        \begin{align*}
            \Expt\bqty{X \pqty{X - 1}} &= \sum_{k = 2}^{n} k\pqty{k - 1} \cdot \frac{n!}{k! \pqty{n - k}!} \cdot p^k \cdot \pqty{1 - p}^{n - k} \\
            &= \pqty{n - 1} n p^2 \sum_{k = 0}^{n - 2} \binom{n - 2}{k} p^k \pqty{1 - p}^{n - 2 - k} \\
            &= \pqty{n - 1} n p^2.
        \end{align*}

        Substituting this back, we have
        \[
            \Var\pqty{X} = \pqty{n - 1} np^2 + np - \pqty{np}^2 = np \pqty{1 - p}.
        \]

        (We will see that using that a binomial random variable is the sum of independent Bernoulli random variables, we may find this using that variance is linear for independent variables.)

        \item Let \(X \sim \Poisson\pqty{\lambda}\) for \(\lambda > 0\), and recall that \(\Expt\bqty{X} = \lambda\).
        
        We have
        \begin{align*}
            \Expt\bqty{X \pqty{X - 1}} &= \sum_{k = 2}^{\infty} k \pqty{k - 1} e^{-\lambda} \frac{\lambda^k}{k!} \\
            &= e^{-\lambda} \lambda^2 \sum_{k = 2}^{\infty} \frac{\lambda^{k - 2}}{\pqty{k - 2}!} \\
            &= \lambda^2
        \end{align*}
        and hence
        \[
            \Var\pqty{X} = \lambda^2 + \lambda - \lambda^2 = \lambda.
        \]
    \end{enumerate}
\end{examples}

\subsection{Covariance}

\begin{definition}[Covariance]
    \label{def:covariance}%
    Let \(X\) and \(Y\) be random variables. We let
    \[
        \Cov\pqty{X, Y} = \Expt\bqty{\pqty{X - \Expt\bqty{X}} \pqty{Y - \Expt\bqty{Y}}}
    \]
    be the \emph{covariance} of \(X\) and \(Y\).
\end{definition}

It is a `measure' of how dependent \(X\) and \(Y\) are.

\begin{proposition}[Properties of Covariance]
    \label{prop:properties-of-covariance}%
    \begin{itemize}
        \item \(\Cov\pqty{X, Y} = \Cov\pqty{Y, X}\). 
        \item \(\Cov\pqty{X, X} = \Var\pqty{X}\).
        \item \(\Cov\pqty{X, Y} = \Expt\bqty{X Y} - \Expt\bqty{X} \Expt\bqty{Y}\).
        \item For \(c \in \RR\), \(\Cov\pqty{cX, Y} = c \Cov\pqty{X, Y}\), and \(\Cov\pqty{c + X, Y} = \Cov\pqty{X, Y}\).
        \item For \(c \in \RR\), \(\Cov\pqty{c, X} = 0\).
    \end{itemize}
\end{proposition}

This gives us a nice formula on the variance of a sum as well.

\begin{proposition}[Variance of Sum]
    \label{prop:variance-of-sum}%
    We have
    \[
        \Var\pqty{X + Y} = \Var\pqty{X} + \Var\pqty{Y} + 2 \Cov\pqty{X, Y}.
    \]    
\end{proposition}

\begin{proof}
    We have
    \begin{align*}
        \Var\pqty{X + Y} &= \Expt\bqty{\pqty{\pqty{X - \Expt\bqty{X}} + \pqty{Y - \Expt\bqty{Y}}}^2} \\
        &= \Var\pqty{X} + 2 \Cov\pqty{X, Y} + \Var\pqty{Y}.
    \end{align*}
\end{proof}

\begin{proposition}[Linearity of Covariance]
    \label{prop:linearity-of-covariance}%
    For random variables \(X, Y\) and \(Z\), we have
    \[
        \Cov\pqty{X + Y, Z} = \Cov\pqty{X, Z} + \Cov\pqty{Y, Z}.
    \]

    More generally, we for any real constants  \(c_1, c_2, \ldots, c_n\) and \(d_1, d_2, \ldots, d_m\) and random variables \(X_1, X_2, \ldots, X_n\) and \(Y_1, Y_2, \ldots, Y_m\), that
    \[
        \Cov\pqty{\sum_{i = 1}^{n} c_i X_i, \sum_{j = 1}^{m} d_j Y_j} = \sum_{i = 1}^{n} \sum_{j = 1}^{m} c_i d_j \Cov\pqty{X_i, Y_j}.
    \]
\end{proposition}

\begin{corollary}[Variance of Sum]
    \label{cor:variance-of-sum}%
    We have
    \[
        \Var\pqty{\sum_{i = 1}^{n} X_i} = \Cov\pqty{\sum_{i = 1}^{n} X_i, \sum_{i = 1}^{n} X_i} = \sum_{i = 1}^{n} \Var\pqty{X_i} + \sum_{i \neq j} \Cov\pqty{X_i, X_j}.
    \]
\end{corollary}

We stated before that the covariance is a `measure' of how independent two random variables are. We first see some important properties on independence that will help with this.

\begin{lemma}[Independence implies Pairwise Independence]
    \label{lem:indep-pairwise-indep}%
    Let \(X_1, X_2, X_3\) be independent, i.e.,
    \[
        \Prob\pqty{X_1 = x_1, X_2 = x_2, X_3 = x_3} = \prod_{i = 1}^{3} \Prob\pqty{X_i = x_i}.
    \]

    Then \(X_1\) is independent of \(X_2\).
\end{lemma}

\begin{proof}
    We have
    \begin{align*}
        \Prob\pqty{X_1 = x_1, X_2 = x_2} &= \sum_{x_3} \Prob\pqty{X_1 = x_1, X_2 = x_2, X_3 = x_3} \\
        &= \sum_{x_3} \Prob\pqty{X_1 = x_1} \Prob\pqty{X_2 = x_2} \Prob\pqty{X_3 = x_3} \\
        &= \Prob\pqty{X_1 = x_1} \Prob\pqty{X_2 = x_2} \sum_{x_3} \Prob\pqty{X_3 = x_3} \\
        &= \Prob\pqty{X_1 = x_1} \Prob\pqty{X_2 = x_2}.
    \end{align*}
\end{proof}

\begin{lemma}[Independence and Expectation of Product]
    \label{lem:independence-expectation-product}%
    Let \(X\) and \(Y\) be independent random variables, and \(\functiondomain{f, g}{\RR}{\RR_{+}}\). Then
    \[
        \Expt\bqty{f\pqty{X} g\pqty{Y}} = \Expt\bqty{f\pqty{X}} \Expt\bqty{g\pqty{Y}}.
    \]
\end{lemma}

\begin{proof}
    Let \(Z = \pqty{X, Y}\), and define
    \[  
        h\pqty{Z} = f\pqty{X} g\pqty{Y}.
    \]

    Then we have
    \begin{align*}
        \Expt\bqty{h\pqty{Z}} &= \sum_{\pqty{x, y}} h\pqty{x, y} \Prob\pqty{Z = \pqty{x, y}} \\
        &= \sum_{\pqty{x, y}} f\pqty{x} g\pqty{y} \Prob\pqty{X = x, Y = y} \\
        &= \sum_{x, y} f\pqty{x} g\pqty{y} \Prob\pqty{X = x} \Prob\pqty{Y = y} \\
        &= \sum_{x} f\pqty{x} \Prob\pqty{X = x} \sum_{y} g\pqty{y} \Prob\pqty{Y = y} \\
        &= \Expt\bqty{f\pqty{X}} \cdot \Expt\bqty{g\pqty{Y}}.
    \end{align*}
\end{proof}

\begin{proposition}[Independence gives Zero Covariance]
    If \(X\) and \(Y\) are independent, then
    \[
        \Cov\pqty{X, Y} = 0.
    \]
\end{proposition}

\begin{proof}
    This follows from \autoref{lem:independence-expectation-product}, since
    \[
        \Expt\bqty{X - \Expt\bqty{X}} = \Expt\bqty{Y - \Expt\bqty{Y}} = 0.
    \]
\end{proof}

The converse does not hold -- in general, for two random variables \(X\) and \(Y\), \(\Cov\pqty{X, Y} = 0\) does \textbf{not} tell us that \(X\) and \(Y\) are independent.

\begin{example}[Zero Variance but Not Independent]
    \label{ex:zero-variance-no-independent}%
    Let \(X_1, X_2, X_3 \sim \Bernoulli\pqty{\frac{1}{2}}\). Set
    \[
        Y_1 = 2 X_1 - 1, \quad Y_2 = 2 X_2 - 1,
    \]
    and
    \[
        Z_1 = Y_1 \cdot X_3, \quad Z_2 = Y_2 \cdot X_3.
    \]

    Then
    \[
        \Expt\bqty{Y_1} = \Expt\bqty{Y_2} = 0.
    \]

    By independence,
    \[
        \Expt\bqty{Z_1} = \Expt\bqty{Y_1 \cdot X_3} = \Expt\bqty{Y_1} \cdot \Expt\bqty{X_3} = 0,
    \]
    and similarly \(\Expt\bqty{Z_2} = 0\).

    Therefore,
    \[
        \Cov\pqty{Z_1, Z_2} = \Expt\bqty{Z_1 \cdot Z_2} = \Expt\bqty{Y_1 \cdot Y_2 \cdot X_3^2} = 0.
    \]

    We want to show \(Z_1\) is not independent of \(Z_2\). Indeed,
    \[
        \Prob\pqty{Z_1 = 0, Z_2 = 0} = \Prob\pqty{X_3 = 0} = \frac{1}{2},
    \]
    but
    \[
        \Prob\pqty{Z_1 = 0} = \Prob\pqty{Z_2 = 0} = \frac{1}{2}.
    \]
\end{example}

Similarly, let \(X_1, \ldots, X_n\) be independent random variables. Then
\[
    \Var\pqty{\sum_{i = 1}^{n} X_i} = \sum_{i = 1}^{n} \Var\pqty{X_i}.
\]

\begin{example}[Binomial as a sum of Bernoulli Random Variables]
    \label{ex:binomial-sum-bernoulli}%
    Let \(S_n \sim \Binomial\pqty{n, p}\). We recall last time that \(\Var\pqty{S_n} = n p \pqty{1 - p}\).

    Now, recall that
    \[
        S_n = X_1 + X_2 + \cdots + X_n
    \]
    where \(X_i\) are independent and identically distributed, with \(X_i \sim \Bernoulli\pqty{p}\).

    Therefore,
    \[
        \Var\pqty{S_n} = \sum_{i = 1}^{n} \Var\pqty{X_i} = np \pqty{1 - p}
    \]
    as desired.
\end{example}